\begin{abstract}
  Vision-based quadruped parkour typically relies on onboard depth cameras, limiting deployment to robots with specialized sensors. We investigate whether a pre-trained \emph{monocular depth estimator} can replace depth cameras for locomotion by producing a canonical depth representation from RGB, without photorealistic rendering or domain randomization.

  We present a two-stage pipeline: a privileged teacher trained via PPO with ground-truth terrain scans, distilled into a visual student via DAgger operating on depth images from Depth~Anything~V2 (DA2). In simulation, this pipeline enables basic parkour on stairs, hurdles, and gaps under conditions characterized by a \emph{Depth Signal Quality} (DSQ) metric. Visual robustness experiments show that DA2-based students tolerate moderate lighting and color perturbations within simulation, though performance degrades on featureless surfaces and under combined perturbations. All results are in simulation; real-world transfer remains future work.
\end{abstract}
